\chapter{Espacios separables y espacios completos}
\label{ch:espacios-separables}
\index{Espacio separable}\index{Conjunto denso}\index{Espacio de Baire}\index{Segunda categoría}

\section{Espacios separables}

La separabilidad es una propiedad de ``tamaño controlable'': un espacio es separable si contiene un subconjunto numerable que aproxima a todos sus puntos. Esta condición, aparentemente técnica, tiene consecuencias profundas: en espacios separables, toda sucesión acotada posee una subsucesión débilmente convergente (teorema de Eberlein--Šmulian); las medidas de probabilidad pueden aproximarse por combinaciones convexas de deltas de Dirac; y muchos espacios de funciones usuales son separables, lo cual facilita su implementación computacional.

\begin{definition}
\label{def:separable}
Un espacio métrico $(X,d)$ es \emph{separable} si posee un subconjunto denso numerable.
\end{definition}

\begin{example}
\label{ej:rn-separable}
$\mathbb{R}^n$ es separable, pues $\mathbb{Q}^n$ es denso y numerable.
\end{example}

\begin{example}
\label{ej:cb-separable}
$C([a,b])$ con la métrica del supremo es separable. Los polinomios con coeficientes racionales forman un conjunto denso numerable (por el teorema de Weierstrass, \cref{teo:weierstrass-aprox}).
\end{example}

\begin{example}
\label{ej:linf-no-separable}
El espacio $\ell^{\infty}$ de las sucesiones acotadas con la métrica del supremo no es separable. La familia de sucesiones de ceros y unos es no numerable y cualquier dos elementos distintos están a distancia $1$.
\end{example}

\section{Conjuntos densos numerables}

\begin{theorem}
\label{teo:base-numerable}
Un espacio métrico es separable si y solo si su topología posee una base numerable.
\end{theorem}
\begin{proof}
Si $\{x_n\}$ es denso, las bolas $B(x_n,q)$ con $q\in\mathbb{Q}^+$ forman una base numerable. Recíprocamente, si $\{U_n\}$ es una base numerable, eligiendo un punto $x_n\in U_n$ en cada elemento no vacío obtenemos un conjunto denso numerable.
\end{proof}

\begin{corollary}
\label{cor:lindelof}
Todo espacio métrico separable es \emph{Lindelöf}: de todo recubrimiento abierto se puede extraer un subrecubrimiento numerable.
\end{corollary}

\section{Espacios completamente métricos}

\begin{definition}
\label{def:completamente-metrico}
Un espacio métrico $(X,d)$ es \emph{completamente métrico} si admite una métrica completa que genere la misma topología. (Nótese que no exigimos que la métrica original sea completa.)
\end{definition}

\begin{theorem}[Alexandrov]
\label{teo:alexandrov}
Un subespacio de un espacio métrico completo es completamente métrico si y solo si es un $G_\delta$ (intersección numerable de abiertos).
\end{theorem}

\section{Espacios de Baire}

Ya definidos en el \cref{ch:grandes-teoremas}, los espacios de Baire merecen un estudio más detallado.

\begin{theorem}[Baire, recapitulación]
\label{teo:baire-recap}
Todo espacio métrico completo es de Baire. La intersección de abiertos densos es densa; ningún espacio de Baire es union numerable de conjuntos nowhere dense.
\end{theorem}

\begin{proposition}
\label{prop:abierto-denso-baire}
En un espacio de Baire, todo abierto no vacío es de segunda categoría en sí mismo (es decir, no es union numerable de nowhere dense).
\end{proposition}

\section{Propiedades fundamentales}

\begin{theorem}
\label{teo:separable-herencia}
Todo subespacio de un espacio métrico separable es separable.
\end{theorem}
\begin{proof}
Si $D$ es denso numerable en $X$ y $Y\subseteq X$, para cada $x\in D$ y $n\in\mathbb{N}$ con $B(x,1/n)\cap Y\neq\emptyset$, elegimos $y_{x,n}\in B(x,1/n)\cap Y$. El conjunto de tales $y_{x,n}$ es denso en $Y$.
\end{proof}

\begin{theorem}
\label{teo:producto-separable}
El producto numerable de espacios métricos separables es separable.
\end{theorem}

\section{Ejemplos importantes}

\begin{example}[Espacio de Hilbert L2]
\label{ej:l2-separable}
El espacio de Lebesgue $L^2([0,1])$ es separable. Los polinomios trigonométricos con coeficientes racionales son densos. En efecto, por el teorema de Stone--Weierstrass (\cref{teo:stone-weierstrass}) aplicado a funciones continuas sobre la circunferencia $[0,2\pi]$ con $0$ y $2\pi$ identificados, el álgebra generada por $\{1,\cos nx,\sin nx:\,n\in\mathbb{N}\}$ es densa en $C([0,2\pi])$ con la métrica del supremo; por densidad, también lo es en $L^2([0,2\pi])$ con la norma cuadrática.
\end{example}

\begin{example}[Espacio de Cantor]
\label{ej:cantor-separable}
El conjunto de Cantor, como subespacio de $\mathbb{R}$, es compacto, perfecto, totalmente disconexo y separable (de hecho, es homeomorfo al producto numerable $\{0,1\}^{\mathbb{N}}$).
\end{example}

\begin{exercise}
Demuestre que todo espacio métrico compacto es separable.
\end{exercise}

\begin{exercise}
Sea $X$ separable y $f\colon X\to Y$ sobreyectiva continua. Demuestre que $Y$ es separable.
\end{exercise}

\begin{exercise}
Pruebe que $\ell^p$ es separable para $1\leq p<\infty$, pero $\ell^{\infty}$ no lo es.
\end{exercise}

\begin{problem}
Un espacio métrico se dice \emph{polaco} si es separable y completamente métrico. Demuestre que $\mathbb{R}\setminus\mathbb{Q}$ es polaco, pero no es completo con la métrica usual.
\end{problem}
